\renewcommand{\columnseprule}{1.5pt}
\begin{multicols*}{2}
\noindent
\rule[0.5\baselineskip]{0.5\textwidth}{1pt}

\noindent
\subsection{Zoomed Straight}
\noindent
1. Input the function $f(x)=\frac{4}{x}$ into your TI-8* and view
it under ZOOM-STANDARD.  Press \Touche[style=function,principal={2ND}]
\Touche[style=function,principal={ZOOM}] and FORMAT your
window to have the GRID ON.  What lattice points (i.e., where
$x$ and $y$ are both integers) does the graph pass through?

\vspace{2cm}
\noindent
2. What do you estimate the rate of change of the function is around 
$f(2)$? Record you guess here.

\vspace{1cm}
\noindent
3.  Using the ZOOM-BOX feature, outline a new window from (0,0) to
approximately (4,4).  Again, draw a ZOOM-BOX, now from about  (1,1)
to (3,3).  Change your \texttt{Xscl} and \texttt{Yscl} to 0.2.  Sketch the 
resultant graph here.

\begin{tikzpicture}[scale=0.5]
    \foreach \x in {0,1,2,3,4,5,6,7,8,9,10}
    \foreach \y in {0,1,2,3,4,5,6,7,8,9,10}
    {
    \draw [gray] (\x,\y) circle (0.1cm);
    }
    \draw [black, -] (-1,-1) -- (-1,11) -- (11,11) -- (11,-1) -- cycle;
  \end{tikzpicture}
  
\noindent
4. Manually zoom in this time, by adjusting the WINDOW so that your
$x$ and $y$ min and max are, respectively, 1.5 and 2.5.  Change both scales
to 0.1.  Press TRACE and move left and right along the graph.  What 
grid points does the equation pass through?  

\vspace{3cm}
\noindent
5.  Make your window 1.99 to 2.01 and your scale 0.002 in both directions.
You should see an almost completely straight line.  TRACE a pair of coordinate,
one just a little left and up from (2,2) and one just a little right and down.  Find
the slope of the line through them, rounding to the nearest integer.  
Record the equation in slope-intercept form,
with that slope, passing through (2,2).

\vspace{3cm}
\noindent
6. Moving the $y_1$ term to the other side, enter your line as another equation
on your TI-8*.  Turn off the GRID and change your scales to 1.  ZOOM-OUT
from the point (2,2) repeatedly, until you approximate ZOOM-STANDARD.
What have you drawn?

\vspace{3cm}
\noindent
7. Using the different quotient, work out the derivative of $f(x)=\frac{4}{x}$ below.
What is $f'(2)$?

\vspace{3cm}
8. It appears that --- on a small enough scale --- inversely proportional graphs
are \textbf{locally linear}.  What other functions do you think might display 
this property?  Experiment on the calculator.  Where you right?

\vspace{4cm}
\noindent
9. Describe what you think the point of this problem set is, using technical vocabulary in complete
sentences.
\end{multicols*}